\section{Introduction}


The Ethereum Virtual Machine (EVM) provides a limited set of precompiled contracts for cryptographic operations: \texttt{ecrecover} (secp256k1 signature recovery), SHA-256, RIPEMD-160, identity, modular exponentiation, and BN254 curve operations (EIP-196/197) \cite{buterin2014ethereum, eip196, eip197}. While these suffice for basic Ethereum functionality, advanced applications---zero-knowledge proofs, aggregate signatures, cross-chain verification, and post-quantum cryptography---require operations that are either absent from the standard precompile set or prohibitively expensive when implemented in Solidity bytecode.

Lux Network addresses this gap through custom precompiles: native Go functions callable from Solidity at designated addresses, executing outside the EVM interpreter at near-native speed. This approach is possible because Lux appchains control their own VM implementations, allowing protocol-level extensions without Ethereum consensus.

This paper makes the following contributions:

\begin{enumerate}[label=(\roman*)]
    \item A comprehensive precompile suite covering five cryptographic families with formal interface specifications.
    \item A gas pricing model based on empirical benchmarking that ensures gas costs reflect wall-clock execution time within a constant factor.
    \item Security analysis of each precompile, including input validation, error handling, and denial-of-service resistance.
    \item Comparison with EIP-2537 (BLS12-381) and other proposed Ethereum precompiles, showing that Lux's implementation achieves lower gas costs through tighter integration.
    \item Post-quantum readiness analysis demonstrating practical ML-DSA (FIPS 204, formerly CRYSTALS-Dilithium) and FALCON verification on-chain.
\end{enumerate}

